\chapter{Geodata and Geospatial Mathematics}

\section{Geographic Coordinate Systems}
Coordinate systems are not only relevant in a theoretical context. Map services, navigation systems, GPS applications, all of them use coordinate systems to orient themselves. Like a coordinate system in 3D mathematics, a geographic coordinate system helps us to define our position on the earth. These systems form the foundation for digital maps for visualizing a geographic environment. This chapter covers the standard coordinate system used in most modern mapping applications. 

\subsection{WGS84}
The World Geodetic System 1984 (WGS84) is the most used geographic reference system in modern times. It is an earth-centered terrestrial reference system, which describes the earth's size, gravity, geomagnetic fields, and shape\ \cite{Web:40}. This system is used as a standard in many applications, including drones, satellites, navigation systems, and digital maps, since it can precisely determine geographical locations\ \cite{Web:41}. It was developed by the National Geospatial-Intelligence Agency, which still maintains it to this day.\ \cite{Web:40}

\begin{figure}[H]
	\centering
	\includegraphics[scale=0.5]{images/geoideb.jpeg}
	\caption{WGS84 system reference,~\cite{Web:42}}\label{fig:WGS84_reference}
\end{figure}	

Figure~\ref{fig:WGS84_reference} above shows how this reference system works. It uses a Cartesian coordinate system with X-, Y-, and Z-axes. The Z-axis goes through the north- and south pole. The center of the system is defined at X=0, Y=0, and Z=0, in the middle of the earth. Points can also be defined in longitude, latitude, and altitude as you can see in the top-right corner of the picture. As defined in the picture, $\phi$ represents latitude, $\lambda$ longitude, and h altitude. They are also defined in the Cartesian system with the Z-X plane being the zero of longitude and the X-Y plane being the zero of latitude.\ \cite{Web:43}

WGS84 provides a reference of the earth's shape using an ellipsoid. Due to it using gravitational-, and geomagnetic fields to provide accurate results, this system needs to be updated since these values change over time and can lead to inaccurate data. These values are updated with the collaboration of organizations including the International Association of Geodesy (IAG) and the International Earth Rotation and Reference Systems Service (IERS).\ \cite{Web:43}



\subsection{Local Tangent Plane}
Although WGS84 being universally applicable, this system is inefficient and impractical for local applications due to the use of angular coordinates (longitude and latitude). For this case, flattening of a small area is more efficient. This method is known as local tangent plane (LTP) which is based on the tangent plane.\ \cite{Web:62}

\begin{quote}
    A tangent plane is a flat surface that touches a curved surface at a single point, representing the local linear approximation of the curved surface at that point.\ \cite{Web:63}
\end{quote}

The tangent plane in LTPs is Earth's rotation axis and the vertical direction of the point. It uses the Cartesian coordinate system. One axis represents the vertical direction, one represents the northern axis, and the last one represents the eastern axis.

\begin{figure}[H]
	\centering
	\includegraphics[scale=0.35]{images/Local-Tangent-Plane-system-of-coordinates-for-GPS-agricultural-applications.png}
	\caption{LTP example,~\cite{Web:64}}\label{fig:LTP_example}
\end{figure}

Figure~\ref{fig:LTP_example} shows an visual example of an LTP\@. $O_{LTP}$ is the origin point of the plane's Cartesian system. This is the point where the LTP touches the ellipsoid shape. The plane is normal to the Z-axis which represents the vertical direction of the point, starting from the Center of mass of the Earth. Usually the Z-axis is called the U-axis. X and Y are represented by the eastern axis E and the northern axis N.

The red point $P_(e,n,z)$ represents a point in space which can now be described in a simple Cartesian coordinate system relative to the local origin of the plane. 

\subsubsection{Conversion of Longitude, Latitude, Altitude to E, N, U coordinates}

Coordinates defined with Longitude, Latitude, and Altitude can be converted into local Cartesian coordinates and back. For the sake of the diploma thesis, only the conversion to LTP coordinates will be covered. 
The first step will be to convert the Longitude, Latitude, and Altitude into earth centred X-, Y-, and Z-coordinates. For this we need to consider a few constant values:

\begin{itemize}
    \item $a = 6378137.0 m$ $\xrightarrow{}$ the semimajor axis which is the equatorial radius
    \item $b = 6356752.3142 m$ $\xrightarrow{}$ the semiminor axis which is the polar radius
    \item $\pi = 3.1415926\dots$ 
\end{itemize}
\noindent\cite{Web:65}

With these values the flatness of the Earth and the eccentricity can be calculated. The eccentricity describes the deviation of an ellipse from a perfect circle, while the flatness of an ellipse is the ratio of the difference between the semimajor and the semiminor axis.

\[
f = \frac{a-b}{a} = 3.3528107\cdot 10^{-3}
\]

\[
e = \sqrt{f(2-f)} = 8.1819191 \cdot 10^{-2}
\]

\noindent\cite{Web:65}

These values are needed to calculate the distance distance from the surface to the z-axis along the ellipsoid normal. This value is represented as a function of latitude $\phi$.

\[
N(\phi) = \frac{a}{\sqrt{1-e^2 \cdot \sin^2(\phi)}} 
\]

With $N(\lambda)$ the X-, Y-, and Z-coordinates can be calculated.

\[
\begin{aligned}
x &= (h+N) \cdot \cos(\phi) \cdot \cos(\lambda)\\
y &= (h+N) \cdot \cos(\phi) \cdot \sin(\lambda)\\
z &= (h + (1-e^2) \cdot N) \cdot \sin(\phi) 
\end{aligned}
\]
\noindent\cite{Web:65}
 
$\lambda$ and $\phi$ represent longitude and latitude. X, Y, and Z are now the position in the Earth-Centered, Earth-Fixed (ECEF) rectangular coordinate system. The next step will be to subtract the coordinates with the origin coordinates of the LTP\@. This moves the origin from the Earth's centre to the LTP origin. 

\[
\begin{bmatrix}
x' \\ y' \\ z' 
\end{bmatrix}
=
\begin{bmatrix}
x \\ y \\ z
\end{bmatrix}
-
\begin{bmatrix}
x_0 \\ y_0 \\ z_0
\end{bmatrix}
\]

$x_0$, $y_0$, and $z_0$ $\xrightarrow{}$ origin coordinates of the LTP \\
$x'$, $y'$, and $z'$ $\xrightarrow{}$ translated origin to the LTP origin \\

\noindent\cite{Web:65}

The final step of the conversion will be to rotate the coordinates. The first rotation aligns the system so that the transformed Y-axis will be parallel to the local East direction.

\[
R_{first} = 
\begin{bmatrix}
-\sin(\lambda) & \cos(\lambda) & 0 \\
\cos(\lambda) & \sin(\lambda) & 0 \\
0 & 0 & 1
\end{bmatrix}
\]

The second rotation rotates the U-axis about the East-axis so it points straight up. This aligns the plane with the LTP normal vector.

\[R_{second} = 
\begin{bmatrix}
1 & 0 & 0 \\
0  & -\sin(\phi) & \cos(\phi) \\
0 & \cos(\phi) & \sin(\phi)
\end{bmatrix}
\]

To finalize the transformation the Earth-centred coordinates need to be multiplied with both rotation matrices. For this calculation the mathematical order of combining transformation needs to be considered to result in the correct coordinates.

\[
\begin{bmatrix}
e \\
n \\
u
\end{bmatrix} = R_{second} \cdot R_{first} \cdot \begin{bmatrix}
x' \\
y' \\
z'
\end{bmatrix}
\]

\noindent\cite{Web:65}



\section{Map Projections}
Map projection allows to visualize the three-dimensional, ellipsoidal Earth on a two-dimensional plane. Due to the Earth's shape, it introduces challenges. The most notable being the distortion of the map. The distortion manipulates the size and shape of objects on Earth. This chapter covers the most popular method and it's web implementation.

\subsection{Mercator \& Web Mercator}
The Mercator projection is the most common choice when it comes to Map projections. This projection maps the Earth on a cylindrical shape while maintaining the right angles and directions\ \cite{Web:46}. It was created by the Flemish mapmaker Gerardus Mercator in 1569\ \cite{Web:46} in a time where world trade had a significant increase in relevance and so did better map projections\ \cite{Web:47}. 

\begin{figure}[H]
	\centering
	\includegraphics[scale=0.5]{images/mercator-projection.jpg}
	\caption{Mercator projection,~\cite{Web:48}}\label{fig:mercator_projection}
\end{figure}

This method proved to be efficient for traveling, since people only had to follow a single straight line and arrive at the exact point on the map without taking the Earth's curvature into consideration. 

Although this method provided the right angles and directions, it was unable to visualize the true size of countries. Since it is projecting an ellipsoid onto a vertical cylinder shape, the texture is heavily distorted at the north and south pole. It is impossible to project a spherical object onto an 2D plane without the help of distortion which means that there is no map projection without it\ \cite{Web:48}. A popular example of distortion would be Greenland which appears to be similar in size compared to the entire continent of Africa:

\begin{figure}[H]
	\centering
	\includegraphics[scale=0.3]{images/greenland1.png}
	\includegraphics[scale=0.3112]{images/greenland2.png}
	\caption{Mercator projection,~\cite{Web:51}}\label{fig:greenland}
\end{figure}


The Mercator projection is still used today in many occasions. With the introduction of digital maps, this method proved to be work well in use with interactive map services\ \cite{Web:49}. Because of this property, the method was implemented in Google Maps and later OpenStreetMap\ \cite{Web:50}. The web implementation is also known as Web Mercator. 


\section{Digital Elevation Models (DEM)}
The Earth is not a perfect ellipsoid. The planet consists of a complex topography with mountains and valleys. For visualizing this unevenness on the surface, digital elevation models (DEMs) come into play. These models provide information about Earth's topography and are mostly created from different sources.\ \cite{Web:52}. 

\begin{figure}[H]
	\centering
	\includegraphics[scale=1.0]{images/dem.png}
	\caption{Digital Elevation Model,~\cite{Web:54}}\label{fig:dem_example}
\end{figure}


There are many ways to create a DEM\@. One method of creating those models is to use synthetic aperture radars. These radars are capable of seeing through clouds as a result of the use of radar waves to collect accurate elevation data during day or night. This method is ideal for frequent cloudy regions.\ \cite{Web:55}

Another way is the ``Light Detection and Ranging'' (LiDAR) method which uses lasers and sensors. The lasers reflect from the surface into the sensors to measure the elevation of the ground. This creates high-resolution and accurate data.\ \cite{Web:55}

Digital elevation model is commonly used as a general term for elevation data. DEMs can be divided into two different types depending on their purpose.\ \cite{Web:56}


\subsection{Digital Terrain Models (DTM)}

Digital terrain models show the surface of the Earth without any vegetation, buildings or other man-made structures. They focus only on the planet's terrain and elevation. Since most elevation data acquirement methods capture all above-ground features of the surface, additional processing has to be done to remove unwanted features above ground.\ \cite{Web:57}

Several filters are applied to the model to remove such structures above ground. In the DSM-to-DTM filtering process, the first step is to remove large structures from the ground. Then pits and small bumps will be removed. After these steps, smoothing filters are applied to restore the natural contours where objects where placed before. The amount of processing needed depends on the quality and structure of the DSM.\ \cite{Web:58}

DTMs are widely used in environmental applications including the planning of structures like roads and canals, and hydrology to plan for example, drainage systems. In terms of map services, a DTM is used as a base layer for geographic information systems. They are used to visualize the bare earth.\cite{Web:59}


\subsection{Digital Surface Models (DSM)}

In contrast to DTMs, which remove all above-ground objects of a DEM, a digital surface model (DSM) shows these features of the Earth's surface. This model provides a detailed and realistic visualization of the Earth with all above-ground features including man-made objects like buildings and bridges alongside the terrain elevation.\ \cite{Web:60}

DSMs have a variety of uses in multiple scenarios. One of the most notable applications is infrastructure planning. In this industry, the models are used to determine land use and simulate the environmental impact due to the expanding infrastructure. In addition to that, these models are also taken advantage of by telecommunications companies for optimizing cell tower placements for optimal performance.\ \cite{Web:61}